\chapter{Análise de Sobrevivência}
\label{cap:sobrevivencia}
% ── índice ──────────────────────────────────────────────────────
\index{análise de sobrevivência|textbf}
\index{Kaplan-Meier, estimador de|textbf}
\index{km@\texttt{km()}|textbf}
\index{Cox, modelo de|textbf}
\index{cox@\texttt{cox()}|textbf}
\index{função de risco|textbf}
\index{hazard ratio|textbf}
\index{censura}
\index{riscos proporcionais, hipótese de}

% ─────────────────────────────────────────────────────────────────────────────
\section{Tempos de falha e censura}

A análise de sobrevivência modela o \emph{tempo até um evento} $T$ —
falência de empresa, morte, abandono escolar. O desafio central é a
\textbf{censura à direita}: quando a observação termina antes do evento,
sabemos apenas $T > c$ (tempo de observação $c$), não o tempo exato.

Dois conceitos fundamentais:

\[
  S(t) = P(T > t)
  \quad\text{(função de sobrevivência)}
  \qquad
  h(t) = \lim_{\delta\to 0} \frac{P(t \leq T < t+\delta \mid T \geq t)}{\delta}
  \quad\text{(função de hazard)}
\]

$S(t)$ decresce de 1 para 0; $h(t)$ mede a taxa instantânea de falha dado
que o indivíduo sobreviveu até $t$.

\subsection*{Kaplan-Meier}

O estimador não-paramétrico de Kaplan-Meier (1958) estima $S(t)$ sem
premissas distribucionais:

\[
  \hat{S}(t) = \prod_{t_j \leq t} \left(1 - \frac{d_j}{n_j}\right)
\]

onde $d_j$ é o número de eventos no tempo $t_j$ e $n_j$ o número em risco
imediatamente antes.

\subsection*{Cox — Hazard Proporcional}

O modelo de Cox (1972) especifica:

\[
  h(t \mid x_i) = h_0(t)\,\exp(\beta'\,x_i)
\]

A função de hazard baseline $h_0(t)$ é deixada não-paramétrica. Os
coeficientes $\beta$ são estimados pela verossimilhança parcial, sem
especificar $h_0$. O exponencial $e^{\beta_j}$ é a \emph{razão de hazard}:
um aumento unitário em $x_j$ multiplica o hazard por $e^{\beta_j}$.

% ─────────────────────────────────────────────────────────────────────────────
\section{Verificação por DGP}

O DGP usa tempos de falha Weibull com parâmetro de forma 1{,}5 e hazard
proporcional em $x \in \{0,1\}$. O verdadeiro log-hazard ratio é $\beta_x = 0{,}8$.

\begin{lstlisting}[caption={DGP de sobrevivência — Weibull com hazard proporcional}]
seed(42)
generate df x      = rbernoulli(0.5)
generate df u      = rnormal()
generate df t_true = exp(-ln(u) / exp(0.8 * x))   // Weibull, beta=0.8
generate df c      = 3.0
generate df time   = t_true * (t_true <= c) + c * (t_true > c)
generate df event  = (t_true <= c)

let m_km  = km(time, event, df)
display m_km

let m_cox = cox(time ~ x, df, event=event)
display m_cox
\end{lstlisting}

\subsection*{Kaplan-Meier}

\begin{lstlisting}[language={}, basicstyle=\ttfamily\small,
  frame=none, numbers=none, backgroundcolor=\color{codbg},
  xleftmargin=0pt, xrightmargin=0pt, breaklines=false]
============== Kaplan-Meier Survival Estimate ==============
Observations:               300
Events:                     246
Median survival:         1.6319
\end{lstlisting}

A mediana de sobrevivência é 1{,}63: metade dos indivíduos experimenta o
evento antes desse tempo.

\subsection*{Cox}

\begin{lstlisting}[language={}, basicstyle=\ttfamily\small,
  frame=none, numbers=none, backgroundcolor=\color{codbg},
  xleftmargin=0pt, xrightmargin=0pt, breaklines=false]
======================= Cox Proportional Hazards Model =======================
Observations: 300   Events: 246   Log-Likelihood: -1235.9213
Concordance:  0.6725
------------------------------------------------------------------------------
Variable  |    coef |  std err |       z |  P>|z| | exp(coef)
x         |  0.7018 |   0.1306 |   5.372 |  0.000 |    2.0174
\end{lstlisting}

$\hat\beta_x = 0{,}702$ (verdadeiro: 0{,}8) — dentro do intervalo de confiança.
O hazard ratio estimado $e^{0{,}702} = 2{,}02$: indivíduos com $x=1$ têm
hazard instantâneo duas vezes maior do que $x=0$. O índice de concordância
$C = 0{,}67$ indica discriminação moderada (0{,}5 = aleatório, 1{,}0 = perfeito).

% ─────────────────────────────────────────────────────────────────────────────
\section{Sintaxe}

\begin{lstlisting}
km(time, event, df)                          // Kaplan-Meier
cox(time ~ x1 + x2, df, event=col_evento)   // Cox PH
\end{lstlisting}

A coluna \lstinline{event} deve ser $1$ para evento e $0$ para censura.

\begin{nota}
  O modelo de Cox assume \emph{hazards proporcionais}: $h(t \mid x=1)/h(t \mid x=0)$
  é constante no tempo. Violações podem ser detectadas via resíduos de
  Schoenfeld ou interações com $\ln(t)$.
\end{nota}
